\begin{abstract}
\noindent
\textbf{Résumé ---} L'intégration de modèles d'intelligence artificielle dans le cadre du projet JUNON et le développement de jumeaux numériques en hydrogéologie soulèvent la question fondamentale de leur fiabilité. Au-delà des métriques de performance classiques, il est crucial de garantir que les prédictions s'appuient sur des facteurs physiques cohérents. Ce travail explore l'apport des méthodes d'explicabilité (XAI) pour répondre à cet enjeu à travers deux axes de recherche principaux.
D'une part, nous étudions l'adaptation des approches d'attribution de caractéristiques aux séries temporelles multivariées afin de permettre aux experts d'évaluer la qualité décisionnelle des modèles. D'autre part, nous évaluons l'apport des explications contrefactuelles pour simuler des scénarios d'action (comme l'impact de la variation des précipitations sur les niveaux piézométriques) et analyser la faisabilité des scénarios générés.
Pour soutenir ces travaux, un environnement applicatif complet a été développé. Il intègre un entrepôt unifié agrégeant diverses données ouvertes (météo, Hub'Eau, piézométrie) et une plateforme web interactive. Cet outil permet la constitution de jeux de données, l'exécution et la comparaison de modèles prédictifs et mécanistiques, ainsi que l'extraction d'explications, posant ainsi les bases des développements futurs menés entre le BRGM et le LIFAT.

\vspace{0.5cm}
\noindent \textbf{Mots-clés :} Intelligence Artificielle Explicable (XAI), Jumeau Numérique, Hydrogéologie, Séries Temporelles Multivariées, Explications Contrefactuelles, Attribution de Caractéristiques.
\end{abstract}
